% Preamble for "The Collatz Problem: Martingales, Momentum and Continued Fractions"
% N. G. Polson and V. Sokolov
% House rules: no bold type anywhere, no dashes in prose, no boxed equations.

\usepackage[paperwidth=7in,paperheight=10in,inner=0.875in,outer=0.625in,
  top=0.85in,bottom=0.85in,headsep=0.25in,footskip=0.4in]{geometry}
\usepackage[T1]{fontenc}
\usepackage[utf8]{inputenc}
\usepackage[sc]{mathpazo}
\renewcommand{\ttdefault}{pcr}
\linespread{1.04}
\usepackage{amsmath,amssymb,amsfonts,mathtools}
\usepackage{amsthm}
\usepackage{graphicx,booktabs,array,caption,subcaption,longtable}
\usepackage[numbers,sort&compress]{natbib}
\usepackage{enumitem}
\usepackage{needspace}
\usepackage[expansion=false,protrusion=true]{microtype}
\usepackage{titlesec,fancyhdr,emptypage}
\usepackage[chapter]{algorithm}
\usepackage{algpseudocode}
\usepackage{listings,xcolor}
\usepackage{imakeidx}
\makeindex[intoc,title=Index,columns=2,options=-s collatz.ist]
\usepackage[hidelinks,bookmarksnumbered]{hyperref}
% Generated by bookrepo.py from book.yaml and results.yaml. Edit those files, not this one.
% \BookDOI is empty until zenodo.py reserve; \BookDOIline prints the whole line, or nothing while it is empty (needs hyperref).
% \BookISBN stays empty here: the free KDP ISBN may appear only in the print edition, and kdp.py interior defines it
% for that build alone, so \BookISBNline prints "Paperback ISBN: ..." in print and nothing in the digital edition.
\newcommand*{\BookVersion}{1.1.0}
\newcommand*{\BookEdition}{1}
\newcommand*{\BookYear}{2026}
\newcommand*{\BookDOI}{10.5281/zenodo.23421852}
\newcommand*{\BookDOIline}{Digital Object Identifier (DOI): \href{https://doi.org/10.5281/zenodo.23421852}{\nolinkurl{10.5281/zenodo.23421852}}}
\providecommand*{\BookISBN}{}
\ifx\BookISBN\empty\newcommand*{\BookISBNline}{}\else\newcommand*{\BookISBNline}{Paperback ISBN: \BookISBN}\fi

\usepackage[capitalise,noabbrev]{cleveref}

\numberwithin{equation}{chapter}
\allowdisplaybreaks
\setlength{\parindent}{1.2em}
\setlist{itemsep=2pt,topsep=4pt}

% ---------------------------------------------------------------- headings
\newcommand{\partimagehook}{}
\titleformat{\part}[display]{\centering\normalfont\Large\scshape}
  {\partname\ \thepart}{14pt}{\LARGE\scshape}[\vspace{18pt}\partimagehook]
\titleformat{\chapter}[display]{\normalfont\Large\scshape\filright}
  {\chaptertitlename\ \thechapter}{10pt}{\LARGE\scshape}
\titlespacing*{\chapter}{0pt}{20pt}{28pt}
\titleformat{\section}{\normalfont\large\scshape}{\thesection}{0.8em}{}
\titleformat{\subsection}{\normalfont\normalsize\itshape}{\thesubsection}{0.8em}{}
\titleformat{\subsubsection}[runin]{\normalfont\normalsize\itshape}{}{0pt}{}[.]
\titlespacing*{\section}{0pt}{16pt plus 4pt minus 2pt}{8pt}
\titlespacing*{\subsection}{0pt}{12pt plus 3pt minus 2pt}{6pt}
\renewcommand{\paragraph}[1]{\par\medskip\noindent\textit{#1}\quad}

% ---------------------------------------------------------------- running heads
\pagestyle{fancy}
\fancyhf{}
\renewcommand{\headrulewidth}{0pt}
\renewcommand{\chaptermark}[1]{\markboth{#1}{}}
\renewcommand{\sectionmark}[1]{\markright{#1}}
\fancyhead[LE]{\small\thepage\hfill\itshape On the Collatz Conjecture}
\fancyhead[RO]{\small\itshape\nouppercase{\leftmark}\hfill\upshape\thepage}
\fancypagestyle{plain}{\fancyhf{}\fancyfoot[C]{\small\thepage}}

% ---------------------------------------------------------------- captions
\captionsetup{font=small,labelfont=sc,labelsep=period,width=0.95\textwidth}
\DeclareCaptionStyle{ruled}{font=small,labelfont=sc,labelsep=period,strut=off}

% ---------------------------------------------------------------- theorems
% amsthm in place of thmtools: thmtools v0.76 raises "Command \c@proposition
% already defined" for shared counters under the TeX Live 2026 kernel.
% Same styles and numbering as the original \declaretheorem block.
\newtheoremstyle{thmstyle}{8pt plus 2pt}{8pt plus 2pt}{\itshape}{0pt}
  {\normalfont\scshape}{.}{0.5em}
  {\thmname{#1}\thmnumber{ #2}\thmnote{ {\normalfont\itshape(#3)}}}
\newtheoremstyle{defstyle}{8pt plus 2pt}{8pt plus 2pt}{\normalfont}{0pt}
  {\normalfont\scshape}{.}{0.5em}
  {\thmname{#1}\thmnumber{ #2}\thmnote{ {\normalfont\itshape(#3)}}}
\theoremstyle{thmstyle}
\newtheorem{theorem}{Theorem}[chapter]
\newtheorem{proposition}[theorem]{Proposition}
\newtheorem{lemma}[theorem]{Lemma}
\newtheorem{corollary}[theorem]{Corollary}
\newtheorem{conjecture}[theorem]{Conjecture}
\theoremstyle{defstyle}
\newtheorem{definition}[theorem]{Definition}
\newtheorem{example}[theorem]{Example}
\newtheorem{remark}[theorem]{Remark}
\newtheorem*{problem*}{Problem}
% \designation{O2}\label{prob:ch11:o2} in an unnumbered problem: \ref prints the designation O2 (\cref would print the enclosing section)
\makeatletter
\newcommand{\designation}[1]{\phantomsection\def\@currentlabel{#1}}
\makeatother
\newtheorem{exercise}[theorem]{Exercise}
\crefname{conjecture}{Conjecture}{Conjectures}
\crefname{exercise}{Exercise}{Exercises}
\renewcommand{\qedsymbol}{$\square$}
\renewcommand{\proofname}{\textit{Proof}}

% ---------------------------------------------------------------- algorithms
\algrenewcommand\algorithmicrequire{\textit{Input:}}
\algrenewcommand\algorithmicensure{\textit{Output:}}
\algrenewcommand\algorithmicfor{\textit{for}}
\algrenewcommand\algorithmicdo{\textit{do}}
\algrenewcommand\algorithmicif{\textit{if}}
\algrenewcommand\algorithmicthen{\textit{then}}
\algrenewcommand\algorithmicelse{\textit{else}}
\algrenewcommand\algorithmicend{\textit{end}}
% line anchors unique across algorithms (numbered by chapter)
\makeatletter
\providecommand*{\theHALG@line}{}
\renewcommand*{\theHALG@line}{\thealgorithm.\arabic{ALG@line}}
% index page ranges are printed with " to " (collatz.ist); link both ends
\let\collatz@hyperpage\hyperpage
\renewcommand*{\hyperpage}[1]{\collatz@hpsplit#1 to \@nil}
\def\collatz@hpsplit#1 to #2\@nil{%
  \if\relax\detokenize{#2}\relax\collatz@hyperpage{#1}%
  \else\collatz@hprange{#1}#2\@nil\fi}
\def\collatz@hprange#1#2 to \@nil{\collatz@hyperpage{#1} to \collatz@hyperpage{#2}}
\makeatother
\algrenewcommand\algorithmicreturn{\textit{return}}
\makeatletter
\renewcommand\fs@ruled{\def\@fs@cfont{\scshape}\let\@fs@capt\floatc@ruled
  \def\@fs@pre{\hrule height.8pt depth0pt \kern2pt}%
  \def\@fs@post{\kern2pt\hrule\relax}%
  \def\@fs@mid{\kern2pt\hrule\kern2pt}\let\@fs@iftopcapt\iftrue}
\makeatother

% ---------------------------------------------------------------- code
\lstset{basicstyle=\ttfamily\footnotesize,keywordstyle=\ttfamily,columns=fullflexible,keepspaces=true,
  frame=lines,framesep=4pt,xleftmargin=6pt,xrightmargin=6pt,
  commentstyle=\itshape\color{black!60},showstringspaces=false,
  breaklines=true,language=Python}

% ---------------------------------------------------------------- macros
\newcommand{\N}{\mathbb{N}}
\newcommand{\Z}{\mathbb{Z}}
\newcommand{\Q}{\mathbb{Q}}
\newcommand{\R}{\mathbb{R}}
\newcommand{\E}{\mathbb{E}}
\newcommand{\Prob}{\mathbb{P}}
\newcommand{\Var}{\operatorname{Var}}
\newcommand{\Cov}{\operatorname{Cov}}
\newcommand{\KL}{D_{\mathrm{KL}}}
\newcommand{\vt}{v_2}
\newcommand{\ds}{\delta^{*}}
\newcommand{\half}{\tfrac{1}{2}}
\newcommand{\Filt}{\mathcal{F}}
\newcommand{\Ztwo}{\mathbb{Z}_2}
\newcommand{\lt}{\log_2 3}
\newcommand{\wt}{\operatorname{wt}}
\newcommand{\off}{\operatorname{off}}
\newcommand{\src}{\operatorname{src}}
\newcommand{\smmdet}{\ensuremath{\mathrm{SMM}_{\mathrm{det}}}}
\newcommand{\smmstate}{\ensuremath{\mathrm{SMM}_{\mathrm{state}}}}

% table of contents without bold
\usepackage{etoolbox}
\makeatletter
\patchcmd{\l@chapter}{\bfseries}{\scshape}{}{}
\patchcmd{\l@part}{\bfseries}{\scshape}{}{}
\makeatother

% no line break hyphenation (house style), with stretch to compensate
\hyphenpenalty=10000
\exhyphenpenalty=10000
\setlength{\emergencystretch}{3em}
